% 
%  chapter5.tex
%  ISELthesis
%  
%  Created by Matilde Pós-de-Mina Pato on 2012/10/09.
%
\chapter{Trabalho Relacionado Estado de Arte}
\label{cha:trab_art}

Nesta secção, serão abordadas algumas  das implementações nas quais inteligência artificial tem estado a ser usada nos últimos anos e algumas sugestões, de outros autores, para o seu uso futuro.

Em seguida, aprofundar-se-á o estudo sobre o uso de inteligência artificial no contexto policial, sobre o qual o trabalho incide, e identificar-se-á algumas necessidades, e que a IA pode auxiliar, que ainda não foram atendidas, como é o caso da automatização do planeamento de investigações criminais.

\section{Inteligência Artificial e as suas aplicações}
\label{sec:ia_aplic}

Desde do momento do lançamento do ChatGPT, em novembro de 2022, que todas as pessoas, empresas e entidades tentam encontrar a melhor forma de integrar esta ferramenta no seu dia-a-dia.

Segundo o livro \textit{Unlocking Artificial Intelligence: From Theory to Applications} [mx - unlock AI], os autores apresentam alguns sistemas que usam Inteligência Artificial, 

Um outro contexto que encontrou na inteligência artificial um bom aliado é o policial. Existem vários processos que podem ter automatizados e novas abordagens que se tornam possíveis com este tipo de tecnologia a integrar o quotidiano das investigação.

\section{Inteligência Artificial no contexto policial}
\label{sec:ia_policia}

\section{Planeamento Automático}
\label{sec:plan_auto}

Um agente de planeamento tem a seu cargo a tarefa de definir uma sequência de ações que permita atingir um objetivo. [mx - IASA Livro].

O planeamento, segundo [mx - autores APL] [mx - APL] envolve a escolha e organização de ações que permitam realizar uma tarefa ou atingir um determinado objetivo. Esta atividade implica, por norma, recorrer a modelos abstratos de uma lista de ações que o agente pode realizar.

\section{LLM e planeamento} %% 
\label{sec:llm_plan}

\section{Agentes e Arquiteturas}
\label{sec:agentes_arq}

Os modelos de base como os LLM, sozinhos não são capazes de resolver sozinhas problemas mais complexos, que envolvem a recolha de informação de diversas fontes, a escolha entre ferramentas e a execução de ações em várias etapas
sob restrições, como latência, permissões, segurança e custo, entre fatores. [mx - art. 29]

Para preencher esta lacuna, existem os agentes de inteligência artificial, que são capazes de captar o estado do ambiente, traçar um plano, aceder a ferramentas, atualizar a memória e verificar os resultados. [mx - art. 29]

\subsection{Conceito de agente e taxonomias}

Os agentes de inteligência artificial são sistemas que combinam modelos de base com raciocínio, planeamento, memória e utilização de ferramentas. [mx - art. 29]. O desenvolvimento de agentes acelerou bastante depois da OpenAI disponibilizar a API para \textit{function-calling}, no final de 2023. [mx - IBM]

Deste modo, no começo de 2026, Bin Xu propõe uma definição formal de \textit{agent transformer}, nome a que o autor atribui a uma abstração de um agente genérico, que procura unificar a terminologia dispersa na literatura: um agente é descrito como um tuplo
\begin{equation}
	\mathcal{A} = (\pi_{\theta}, \mathcal{M}, \mathcal{T}, \mathcal{V}, \mathcal{E}),
	\label{eq:agent_transformer}
\end{equation}
em que \(\pi_{\theta}\) representa a  \textit{transformer} (LLM/VLM), \(\mathcal{M}\) a memória, \(\mathcal{T}\) o conjunto de ferramentas, \(\mathcal{V}\) os verificadores/críticos e \(\mathcal{E}\) o ambiente [mx - art. 29]. Esta formalização é acompanhada de um \textit{loop} de execução explícito. 

Primeiro dá-se a \textbf{fase de observação} com a recolha de informação do ambiente. Depois a seguem-se a fase de \textbf{recuperação de memória relevante}, a de \textbf{proposta de ação} que recorre à política \(\pi_{\theta}\) condicionada pela observação e memória recuperada. Por fim, dá-se a \textbf{fase de validação} com \(\mathcal{V}\), a de \textbf{execução} que executa a ferramenta escolhida pela ação, que acaba no final por atualizar o ambiente e a memória. Este loop é apresentado nas seguintes equações?

\begin{align}
	\mathcal{A} &= (\pi_{\theta}, \mathcal{M}, \mathcal{T}, \mathcal{V}, \mathcal{E}), \tag{1} \\[4pt]
	o_t &\leftarrow \mathrm{Obs}(\mathcal{E}_t), \quad m_t \leftarrow \mathrm{Retrieve}(\mathcal{M}_t, o_t), \tag{2} \\[4pt]
	\tilde{a}_t &\sim \pi_{\theta}(\cdot \mid o_t, m_t), \quad \hat{a}_t \leftarrow \mathrm{Validate}(\mathcal{V}, \tilde{a}_t), \tag{3} \\[4pt]
	\mathcal{E}_{t+1} &\leftarrow \mathrm{Exec}(\mathcal{E}_t, \mathcal{T}, \hat{a}_t), \quad \mathcal{M}_{t+1} \leftarrow \mathrm{Update}(\mathcal{M}_t, o_t, \hat{a}_t, \mathcal{E}_{t+1}). \tag{4}
\end{align}

Neste contexto, o autor coloca o agente como um \textit{budgeted controller} e ciente do risco: ações de baixo risco correm com quase nenhuma intervenção, enquanto ações de alto impacto exigem verificação adicional ou confirmação humana [mx - art. 29].

%%%%%%%%%%%%%%%%
%% TAXONOMIA ART. 22 / ART . 23)
%%%%%%%%%%%%%%%

No mesmo artigo [mx - art. 29], o autor apresenta uma taxonomia que se aplica ao \textbf{domínio e aplicação}. Esta categorização organiza os agentes pelo \emph{tipo de tarefa e de ambiente} que enfrentam e leva em conta três eixos: (i) o \textbf{foco dominante da interação} (texto/ferramentas, incorporação física, ambientes simulados), (ii) o \textbf{alvo generativo} (conteúdo, mundos, experiências) e (iii) o \textbf{substrato de raciocínio} (conhecimento, lógica, emoção, estrutura neuro-simbólica). A partir destes eixos, surgem seis categorias: \textbf{agentes generalistas}, \textbf{agentes incorporados}, \textbf{agentes de simulação e ambientes}, \textbf{agentes generativos}, \textbf{agentes de conhecimento e inferência lógica}, e \textbf{agentes LLM/VLM}. O autor caracteriza esta taxonomia como \textit{pragmática}, pois pretende refletir as restrições que mais condicionam o design como a observabilidade, segurança, latência, verificação e avaliação. [mx - art. 29]

\begin{itemize}
	\item \textbf{Agentes generalistas:} resolvem tarefas heterogéneas (código, navegação, análises, \textit{workflows} empresariais) com um núcleo de política partilhado, ferramentas modulares e memória.
	\item \textbf{Agentes incorporados:} atuam no mundo físico (robótica, dispositivos inteligentes), onde as ações têm custos reais e a segurança é primordial.
	\item \textbf{Agentes de simulação e ambientes:} operam em ambientes virtuais (jogos, \textit{web sandboxes}, mundos sintéticos), onde a interação é mais barata e escalável.
	\item \textbf{Agentes generativos:} produzem artefactos ou experiências além de uma resposta única (histórias, cenas, personagens orientadas pelo diálogo, simulações sociais).
	\item \textbf{Agentes de conhecimento e inferência lógica:} priorizam a correção sob restrições explícitas, recuperando evidências, aplicando regras e produzindo traços auditáveis.
	\item \textbf{Agentes LLM/VLM:} têm a capacidade dirigida principalmente pelo backbone do próprio modelo LLM/VLM utilizado.
\end{itemize}



\subsection{Padrões de arquitetura}
No documento \textit{Building Effective AI Agents: Architecture Patterns and Implementation Frameworks} [mx - Anthropic Doc], a Anthropic apresenta alguns casos de utilização e padrões de arquitetura de agentes. 

No \textit{Single Agent}, um agente percebe o ambiente, decide os próximos passos e age para atingir um objetivo, tudo isto em \textit{loop}. Esta abordagem revelou-se boa para resolver problemas gerais num domínio específico; porém não garante uma precisão de 100\% na primeira resposta.

No \textit{Multi-Agents}, é apresentado um padrão que permite coordenar vários agentes especializados para resolver problemas complexos. Comparativamente ao padrão \textit{Single Agent}, este consome mais \textit{tokens}.

Este padrão pode seguir duas abordagens: \textbf{Hierárquica/Supervisora} e \textbf{Colaborativa}. Na primeira, existe um agente "supervisor" que delega as tarefas a subagentes especializados. Porém, apresenta-se a necessidade de gerir o contexto, porque o supervisor pode ficar sobrecarregado. Já na segunda abordagem, existem vários agentes autónomos que falam entre si, negociam papéis a desempenhar e colaboram sem a necessidade de um supervisor. Contudo, podem surgir problemas na comunicação, o que pode levar a comportamentos imprevisíveis.

O terceiro padrão, \textit{Agentes de Fluxo de Trabalho} refere-se as estruturas pré-definidas que orquestram como os agentes operam, comunicam e transferem tarefas entre si. Estão estruturas podem pertencer a um de três tipos: \textbf{Sequencial}, \textbf{Avaliador-Otimizador} ou \textbf{Paralelo}. No primeiro tipo, as tarefas são feitas numa ordem pré-definida, onde cada agente na sequência faz o seu trabalho e passa o resultado para o seguinte na linha. No segundo tipo, tem-se um agente que gera conteúdo e um outro que o avalia e dá \textit{feedback}. Este ciclo repete-se até que a qualidade desejado seja atingida.

Já no último tipo deste padrão apresentado pela Anthropic, encontram-se vários agentes a executarem tarefas de forma independente e simultaneamente. Os resultados produzidos por cada agente são agregados ou sintetizados no final.

Em 2024, um grupo de investigadores propôs uma taxonomia com cinco grandes categorias para classificar os vários métodos que recorrem a LLM como módulos de planeamento. [mx - art. 22].

As categorias propostas são: \textbf{Decomposição de Tarefas}, \textbf{Seleção Multi-plano}, \textbf{Planeamento Auxiliado por Planeador Externo}, \textbf{Reflexão e Refinamento} e \textbf{Planejamento Aumentado por Memória}.

Na Decomposição de Tarefas, os autores apresentam a filosofia de \textit{dividir para conquistar} como ideia principal ao decompor uma tarefa complexa em subtarefas mais simples e planeá-las sequencialmente. Os métodos para esta decomposição podem ser do tipo \emph{decomposition-first}, no qual o processo ocorre antes do planeamento das subtarefas, ou \emph{interleaved}, onde a decomposição e o planeamento das subtarefas são intercalados.

Na categoria de Seleção Multi-plano procura-se compensar a incerteza e a possível falta de otimização de um único plano gerado pelo LLM, ao gerar vários planos candidatos e, no fim, selecionar o mais adequado. A geração pode recorrer a amostragem ou a mecanismos de proposta explícita, enquanto a seleção utiliza estratégias como votação maioritária, procura em árvore ou A*.

No Planeamento Auxiliado por Planeador Externo o LLM é usado sobretudo para formalizar a tarefa e o ambiente, depois delega a geração do plano num planeador externo, que pode ser simbólico, como PDDL, ASP ou Fast-Downward, ou neuronal, como modelos treinados com aprendizagem por reforço ou \textit{imitation learning}. O objetivo é tirar partido da estabilidade e interpretabilidade dos planeadores externos para lidar com restrições complexas e melhorar a eficiência e viabilidade dos planos. O \textbf{LLLM+P} encontra-se entre os inclui o \textbf{LLM+P}.

A Reflexão e Refinamento é uma categoria que enfatiza a melhoria iterativa do plano através de ciclos de geração com \textit{feedback}, reflexão e refinamento. O agente analisa falhas ou o \textit{feedback} do ambiente, gera reflexões em linguagem natural e ajusta o plano em tentativas posteriores, o que aumenta a tolerância a erros e ajuda a evitar ciclos de pensamento.

Por último, na categoria de Planeamento Aumentado por Memória enriquece-se o planeamento com um módulo de memória que guarda conhecimento, experiências passadas ou informação específica do domínio e que recupera a relevante no momento do planeamento. A memória pode assumir a forma de RAG, com armazenamento externo e recuperação por similaridade, ou de memória incorporada, através de \textit{fine-tuning} do LLM com amostras experimentais do agente.

\subsection{\textit{Model Protocol Context}}
\label{sec:mcp_arc}